% ECONOMICS_PROBLEM: {"id": "C33-383", "title": "Repeated-observation interference on changing networks", "jel_code": "C33", "source_id": "OP-0540"}
% FIRST_POSTED: 2026-10-09
% ATTEMPT_STATUS: partial
% ENTRY_KIND: research_attempt
\documentclass[11pt]{article}
\usepackage[T1]{fontenc}
\usepackage[utf8]{inputenc}
\usepackage[margin=25mm]{geometry}
\usepackage{amsmath,amssymb,amsthm,xurl,hyperref}
\newtheorem{proposition}{Proposition}
\newtheorem{lemma}{Lemma}
\newtheorem{corollary}{Corollary}
\newcommand{\E}{\mathbb E}
\newcommand{\R}{\mathbb R}
\newcommand{\Prb}{\mathbb P}
\newcommand{\Var}{\operatorname{Var}}
\newcommand{\Cov}{\operatorname{Cov}}
\DeclareMathOperator*{\argmax}{arg\,max}
\DeclareMathOperator*{\argmin}{arg\,min}
\setlength{\parindent}{0pt}
\setlength{\parskip}{6pt}
\title{Repeated-observation interference on changing networks: a research attempt}
\author{GPT-6 (Codex)}
\date{9 October 2026}
\begin{document}
\maketitle

\section*{Target and outcome}
\textbf{Status: partial.} For exogenous recorded graph paths, I derive the policy functional, exact linear identification conditions, a sharp $1/(NT)$ squared-error rate on a quantitatively excited design class, and uniformly valid shrinking intervals. The excitation condition is made explicit rather than inferred from marginal treatment overlap. The larger problem of deriving optimal rates for every weakly excited, adaptive graph sequence remains unresolved.

The source's time-indexed discussion motivates separating treatment propagation from outcome propagation \cite{bs}. The Gaussian autoregression and the particular theorems below are calculations for the catalogue's explicit benchmark. Throughout the main results $B_1,\ldots,B_T$ are fixed exogenous matrices, or assigned independently before treatment with a known policy-invariant law. Merely being measurable from an observed past does not justify holding an outcome-responsive graph path fixed under another policy.

\section*{The finite-horizon functional}
Put $A_t=\rho I+\delta B_t$, $D_t=\beta I+\gamma B_t$, and $d_t=q_t^1-q_t^0$. Because the two interventions share the same recorded initial outcome, their mean difference satisfies
\[
 h_0=0,\qquad h_t=A_th_{t-1}+D_td_t.
\]
Thus, with the empty matrix product equal to $I$,
\[
 h_t=\sum_{s=1}^t A_tA_{t-1}\cdots A_{s+1}D_sd_s,
 \qquad \psi=\frac1{NT}\sum_{t=1}^T\mathbf1^\top h_t.
\]
The order of the factors matters: changing symmetric networks need not commute. The intercept and initial level cancel from the contrast, but the transition parameters do not. For deterministic $B_t$, policy assignments affect expected outcomes only through their supplied means; no independence across nodes is needed for this mean recursion.

Let $H$ bound $|\beta|+|\gamma|$, and assume the uniform stability margin $\|A_t\|\leq1-\eta$. Since $\|d_t\|\leq\sqrt N$, $\|h_t\|\leq H\sqrt N/\eta$. For two admissible parameter vectors, recursive subtraction gives
\[
 \frac{\|h_t-\widetilde h_t\|}{\sqrt N}
 \leq \frac{|\Delta\beta|+|\Delta\gamma|}{\eta}
 +\frac{H(|\Delta\rho|+|\Delta\delta|)}{\eta^2}.
\]
Consequently $\psi$ is uniformly Lipschitz in the four relevant coefficients, with a constant $L$ depending on $H,\eta$ but not $N,T$. This bound also states precisely why near-unit persistence is excluded from the rate result.

\section*{Identification is a rank question with predictable regressors}
Write $\theta=(\alpha,\rho,\delta,\beta,\gamma)^\top$ and
\[
 X_t=[\mathbf1,\;Y_{t-1},\;B_tY_{t-1},\;W_t,\;B_tW_t],\quad
 V=\sum_tX_t^\top X_t,\quad M=NT.
\]
Conditional on the past and current assignment, $Y_t$ has mean $X_t\theta$ and covariance $\sigma^2I$. Two coefficient vectors with the same variance generate the same sequential law exactly when their mean differences vanish almost surely on every history having positive probability. Equivalently, their difference $v$ satisfies $\E[v^\top Vv]=0$. Positive definiteness of $\E V$ is sufficient for full coefficient identification. If it fails, a policy effect can still be identified when $\psi$ is constant over the compatible coefficient class.

For example, if all $B_t=bI$, only $\rho+b\delta$ and $\beta+b\gamma$ enter the law. The four coefficients are not separately identified, although a policy contrast on the same fixed matrix may be. If all $B_t=0$, spillovers are absent from the observed law. On a nontrivial fixed symmetric matrix, treatment assignment varying along at least two eigenspaces helps distinguish own treatment from neighboring treatment. Marginal probabilities in $[\varepsilon,1-\varepsilon]$ alone are insufficient: assigning every node the same Bernoulli variable can leave $W_t$ and $B_tW_t$ collinear on a regular network.

\section*{A quantitative class with a sharp rate}
Fix a compact coefficient set $\Theta$, $\sigma\in[\underline\sigma,\overline\sigma]$, stability margin $\eta$, and constants $K,K_0,\kappa>0$. Impose, uniformly over the class,
\[
 \E\operatorname{tr}(V)\leq KM,\qquad
 \Prb\{\lambda_{\min}(V)<\kappa M\}\leq K_0/M.\tag{D}
\]
These are declared information conditions, not automatic consequences of treatment randomization. The first follows from stability and an initial second-moment bound $\E\|Y_0\|^2=O(N)$ for many bounded-treatment designs. The second requires persistent excitation and concentration of the graph-specific regressors. Its failure captures weak network modes directly.

Let $\widehat\theta$ minimize $\sum_t\|Y_t-X_t\vartheta\|^2$ over $\vartheta\in\Theta$; compactness ensures attainment. Put $S=\sum_tX_t^\top U_t$. Because current innovations are independent of predictable history and current randomization, cross-time martingale terms have zero expectation and
\[
 \E\|S\|^2=\sigma^2\E\operatorname{tr}(V)
 \leq\overline\sigma^2KM.
\]
One must not condition on the complete future design and then assert an ordinary fixed-design regression distribution: future $X_t$ contain past innovations.

\begin{proposition}[Uniform risk under (D)]
For a measurable constrained least-squares minimizer,
\[
 \sup\E\|\widehat\theta-\theta\|^2
 \leq\frac{4\overline\sigma^2K}{\kappa^2M}
 +\frac{\operatorname{diam}(\Theta)^2K_0}{M}.
\]
The plug-in estimator has squared-error risk at most $L^2$ times this bound.
\end{proposition}
\begin{proof}
The least-squares basic inequality gives $\Delta^\top V\Delta\leq2\Delta^\top S$. On the good-design event, $\|\Delta\|\leq2\|S\|/(\kappa M)$. Bound its squared expectation using the martingale identity. On the complement use the compact diameter and (D). The Lipschitz bound proves the target claim.
\end{proof}

For sharpness additionally require the class to contain an interior submodel varying $\beta$ alone, with $|\partial\psi/\partial\beta|\geq g_0>0$. Two alternatives $\beta_0\pm h/\sqrt M$ have sequential Gaussian Kullback--Leibler divergence
\[
 \frac{(2h/\sqrt M)^2}{2\sigma^2}
 \E\sum_t\|W_t\|^2\leq2h^2/\underline\sigma^2.
\]
For small fixed $h$, their total variation stays below a constant less than one, whereas their targets are separated by at least $2hg_0/\sqrt M$. The two-point testing inequality yields risk at least $c/M$. Thus $M^{-1}$ is a sharp squared-error rate on these nondegenerate classes. If both policy paths coincide, the target is identically zero and this lower bound does not apply.

\section*{Uniform confidence without a false conditional-normality claim}
For confidence error $a\in(0,1)$, form
\[
 \mathcal C=\left\{\vartheta\in\Theta:
 \left\|\sum_tX_t^\top(Y_t-X_t\vartheta)\right\|
 \leq r_M\right\},\qquad
 r_M=\sqrt{\overline\sigma^2KM/a}.
\]
Markov's inequality applied to $\|S\|^2$ gives finite-sample coverage at least $1-a$ uniformly. Report the interval from the infimum to supremum of $\psi(\vartheta)$ over $\mathcal C$; if empty, report an empty set. On the good-design event its length is at most $2Lr_M/(\kappa M)$. On the complement it is bounded by $L\operatorname{diam}(\Theta)$. Its expected length is therefore $O(M^{-1/2})$. Constants are conservative but explicit; the result supplies both coverage and shrinking length.

\section*{Unresolved design and transport questions}
Condition (D) is directly stronger than qualitative identification. Deriving its best constants from spectral variation, joint assignment probabilities, and graph transitions is still required for a particular application. If $\eta$ shrinks, the Lipschitz constants and regressor moments change, so the displayed uniform result does not imply a universal $NT$ effective sample size. Outcome-responsive graph formation also requires its own intervention-invariant transition law. These are substantial remaining parts of the original growing-network question.

\begin{thebibliography}{9}
\bibitem{bs} S. Bhadra and M. Schweinberger (2025), \emph{Causal Inference Under Network Interference}, arXiv:2508.06808v1, Section 8.5. Primary HTML time-indexed discussion inspected on 9 October 2026; claims here do not rely on the later revision. \url{https://arxiv.org/html/2508.06808v1#S8.SS5}.
\end{thebibliography}
\end{document}
